
\documentclass[11pt]{article}

\usepackage[T1]{fontenc}
\usepackage[utf8]{inputenc}
\usepackage{amsmath,amssymb,amsfonts}
\usepackage{graphicx}
\usepackage{bm}
\usepackage{siunitx}
\usepackage{geometry}
\usepackage[colorlinks=true,linkcolor=black,citecolor=blue,urlcolor=blue]{hyperref}
\usepackage[noabbrev,nameinlink]{cleveref}
\geometry{margin=1in}

\title{Self-starting Siphon Priming: Geometry, Dynamics, and Loss}
\author{Stelios Raptopoulos}
\date{18 November 2025}

\begin{document}
\maketitle

\begin{abstract}
Self-starting siphons use geometry to create a momentum-driven water slug that climbs
over a crest without external priming. This article combines a qualitative ``mental model'' of
the priming process with a one-dimensional slug-flow model that tracks the advancing front
$L(t)$ and the container level $h_c(t)$, including frictional and minor losses and the transition
to a draining phase once the pipe is full.
\end{abstract}

\tableofcontents

\bigskip

\section{Introduction}

Self-starting siphons differ from conventional inverted-U siphons in that they do not require
an externally imposed pressure difference or pre-filling of the tube. Instead, the geometry is
arranged so that an initial drop converts gravitational potential energy into kinetic energy in
a finite slug of water, and the resulting momentum carries the slug over a crest before friction
and gravity bring it to rest. Once the crest has been crossed, the device behaves like a standard
siphon and continues to drain until the head between the reservoir and outlet is exhausted.

The aim of this article is to connect this qualitative picture to a one-dimensional slug-flow
model that is simple enough for analysis and simulation but rich enough to capture the
essential dynamics. We introduce the centreline geometry, the arclength coordinate $L(t)$
and a volume-based description of the container level, derive the basic lossless equation
of motion for the filling phase, add distributed and local (minor) losses, and then write
down the coupled system that governs the draining phase once the pipe is full.
At the level of this slug model, the combined effect of pressure and gravity on the water in
the pipe can be summarised very compactly: the net driving force is
\[
  F_{\text{drive}} = \rho g A\bigl[h_c(t) - h(L(t))\bigr],
\]
where $h_c(t)$ is the container free-surface elevation and $h(L(t))$ is the centreline height
of the water--air front. This driving term is non-zero whenever the front is below the free
surface, i.e.\ while $h(L(t))<h_c(t)$, so after the initial drop into the lower leg the slug
generally re-crosses the container level with finite speed before decelerating towards the
crest. By contrast, in a simple inverted-U siphon the driving head vanishes at the crest
after the initial drop, so the flow does not enjoy the same momentum advantage.

In practice there are two experimental ``wetting'' procedures (pre-filling the leg up to a point
near $B$, or dipping the inlet until the interface reaches $B$) which both lead, at the start
of the dynamical filling phase, to the same initial conditions: the front sits at the end of
segment $AB$ and is momentarily at rest, while the container level passes through $A$.
We take $t=0$ at that moment and show below that both procedures collapse to
\[
  t_0 = 0,\qquad L(0) = L_1,\qquad \dot L(0) = 0,\qquad h_c(0) = h_0.
\]

Self-starting siphon behaviour and priming have been discussed in the fluid-mechanics and
physics-education literature; see, for example, Potter and Barnes~\cite{potter_barnes}, Jumper and Stanchev~\cite{jumper_stanchev},
Zheng et al.~\cite{zheng_siphon}, Ortloff and Kassinos~\cite{ortloff_kassinos}, and Sedlar et al.~\cite{sedlar_siphon}.
The slug-flow model developed below focuses on the transient filling and draining dynamics
of a self-starting siphon, including frictional and minor losses.

Throughout, we make the following modelling assumptions:
\begin{itemize}
  \item The fluid is incompressible and Newtonian with constant density $\rho$ and viscosity
        $\mu$ (e.g.\ water at room temperature).
  \item The pipe is rigid, stationary, and has a uniform circular cross-section of diameter $D$
        and area $A$; the centreline geometry may be curved but lies in a vertical plane.
  \item The container is large enough that the velocity at the free surface is negligible, so the
        container remains hydrostatic at each instant.
  \item Both the container free surface and the outlet discharge to the same atmospheric pressure,
        and the external air dynamics are neglected.
\end{itemize}

\section{Geometry and coordinate systems}
\label{sec:geometry}

We model a self-starting siphon consisting of a rigid container connected to a rigid pipe.
Water flows from the container into the pipe under gravity. The geometry is fully described
in a two-dimensional vertical plane by five node points and two coordinate systems.

\subsection{Global Cartesian coordinate system $(x,h)$}

We introduce a global Cartesian system:
\begin{itemize}
  \item Horizontal coordinate: $x$ (in metres), increasing to the right.
  \item Vertical coordinate: $h$ (in metres), increasing upward.
\end{itemize}
We choose the vertical origin such that $h=0$ at the container base, so all heights (including $h_0,h_1,\dots,h_4$) are measured vertically upward from this level. In particular, $h_0$ is the initial water depth in the container.
All heights in the model (container level, pipe elevations) are measured as values of $h$.
The pipe centreline lies entirely in this $(x,h)$ plane.

\subsection{Pipe centreline: five node points and four segments}

The pipe centreline is defined by five node points
\[
A(0,h_0),\quad B(x_1,h_1),\quad C(x_2,h_2),\quad D(x_3,h_3),\quad E(x_4,h_4),
\]
connected by straight segments:
\begin{align*}
  &\text{Segment 1: from $A$ to $B$, length} &&
    L_1 = \sqrt{(x_1 - 0)^2 + (h_1 - h_0)^2},\\
  &\text{Segment 2: from $B$ to $C$, length} &&
    L_2 = \sqrt{(x_2 - x_1)^2 + (h_2 - h_1)^2},\\
  &\text{Segment 3: from $C$ to $D$, length} &&
    L_3 = \sqrt{(x_3 - x_2)^2 + (h_3 - h_2)^2},\\
  &\text{Segment 4: from $D$ to $E$, length} &&
    L_4 = \sqrt{(x_4 - x_3)^2 + (h_4 - h_3)^2}.
\end{align*}
The total pipe length is
\[
  L_{\max} = L_1 + L_2 + L_3 + L_4.
\]
The inlet at the container wall is at node $A$; the outlet is at node $E$.

We choose the vertical datum such that at $t=0$ the container free surface passes through
node $A$. Thus
\begin{itemize}
  \item the vertical coordinate of $A$ is $h_0$,
  \item and the initial container level is $h_c(0) = h_0$.
\end{itemize}
The outlet elevation is
\[
  h_{\text{out}} = h(L_{\max}) = h_4.
\]

\subsection{Centreline elevation as a function of arclength, $h(L)$}

We use the same symbol $L$ both for:
\begin{itemize}
  \item a generic arclength coordinate along the pipe centreline (geometry), and
  \item the wetted length $L(t)$ (state variable) during the filling phase.
\end{itemize}
When we write $L$ without an explicit time argument, we mean a generic arclength coordinate
measured from node $A$; when we write $L(t)$, we mean the current wetted length.

As a geometric coordinate, $L$ runs from the inlet to the outlet:
\[
  L = 0 \text{ at } A,\quad
  L = L_1 \text{ at } B,\quad
  L = L_1 + L_2 \text{ at } C,\quad
  L = L_1 + L_2 + L_3 \text{ at } D,\quad
  L = L_{\max} \text{ at } E.
\]

The elevation of the pipe centreline as a function of arclength is denoted $h(L)$. Since
each segment is straight, $h(L)$ is piecewise linear in $L$:
\[
h(L) = \begin{cases}
  h_0 + \dfrac{L}{L_1}(h_1 - h_0), & 0 \le L \le L_1 \text{ (segment $A$--$B$)}, \\[1em]
  h_1 + \dfrac{L - L_1}{L_2}(h_2 - h_1), & L_1 \le L \le L_1 + L_2 \text{ (segment $B$--$C$)}, \\[1em]
  h_2 + \dfrac{L - L_1 - L_2}{L_3}(h_3 - h_2), & L_1 + L_2 \le L \le L_1 + L_2 + L_3 \text{ (segment $C$--$D$)}, \\[1em]
  h_3 + \dfrac{L - L_1 - L_2 - L_3}{L_4}(h_4 - h_3), & L_1 + L_2 + L_3 \le L \le L_{\max} \text{ (segment $D$--$E$)}.
\end{cases}
\]
This function $h(L)$ encodes the entire geometry of the pipe (drops, rises, and crest).
For later analytic work it is convenient to introduce cumulative lengths
\[
  S_0 = 0,\quad S_1 = L_1,\quad S_2 = L_1 + L_2,\quad S_3 = L_1 + L_2 + L_3,\quad S_4 = L_{\max},
\]
and to rewrite $h(L)$ on each segment in affine form
\begin{equation}
  h(L) = c_i + s_i L,\qquad
  s_i = \frac{h_i - h_{i-1}}{L_i},\qquad
  c_i = h_{i-1} - s_i S_{i-1},
  \label{eq:linear-hi}
\end{equation}
for $L\in [S_{i-1},S_i]$. If segment $i$ makes an angle $\theta_i$ with the horizontal then
its slope $s_i$ is simply $s_i = \sin\theta_i$, so only the vertical projection of each leg
enters the priming analysis.

\subsection{Wetted length coordinate $L(t)$ and the advancing front}

During the filling phase the pipe is not yet completely full. We introduce the wetted length:
\begin{itemize}
  \item $L(t)$ is the length of pipe actually filled with water at time $t$,
  \item measured along the centreline from node $A$,
  \item with $0 < L(t) \le L_{\max}$.
\end{itemize}
The water--air interface (front) is located at the arclength position $L(t)$, and its elevation is
\[
  h_{\text{front}}(t) = h(L(t)).
\]
When the front has advanced all the way to the outlet, $L(t) = L_{\max}$ and the pipe becomes
completely full.

\section{Container and pipe properties}
\label{sec:container}

We consider a rigid vertical cylindrical container of inner radius $R_c$. Its geometric
cross-sectional area is
\[
  A_{\text{cyl}} = \pi R_c^2.
\]
In the model we do not work directly with $A_{\text{cyl}}$, but with an effective free-surface
area $A_0$ defined as the cross-sectional area actually occupied by water in the container,
after subtracting any internal solid obstructions (such as the pipe wall). In other words,
$A_0$ is the net area of the water region in a horizontal cross-section, not including the
pipe interior.

\begin{itemize}
  \item \emph{Container free-surface area (assumed constant):}
  \[
    A_0\ \text{[m$^2$]}.
  \]
  $A_0$ is the effective area of the water region in the cylindrical container, i.e.\ the
  cylindrical cross-section minus the cross-sectional area of any internal obstructions such
  as the pipe wall. The interior of the pipe (whether filled with water or air) is not part of
  the container control volume. With this definition, the water volume in the container is
  \[
    V_{\text{cont}}(t) = A_0 h_c(t),
  \]
  where $h_c(t)$ is the container free-surface elevation measured against the cylinder wall.
  \item \emph{Pipe cross-sectional area:}
  \[
    A\ \text{[m$^2$]},
  \]
  with circular diameter $D$. The water volume inside the pipe at time $t$ is
  \[
    V_{\text{pipe}}(t) = A L(t).
  \]
\end{itemize}

By definition, any portion of the pipe interior, even if it lies below the current free surface
and is geometrically inside the cylindrical shell, is counted as part of the pipe volume
$V_{\text{pipe}}$, not the container volume $V_{\text{cont}}$. The displacement effect of the
pipe wall on the free surface is already built into the choice of $A_0$ as an effective area,
so $V_{\text{cont}} = A_0 h_c$ and $V_{\text{pipe}} = A L(t)$ together account for all the
water in the system without double-counting.

\subsection{On the effective area $A_0$ for complex internal pipe layouts}

The physical prototype motivating this model has a meandering pipe inside the container: a
horizontal plane at a given elevation may intersect zero, one, or several legs of the pipe.
In such a configuration, the area occupied by water in a horizontal cross-section is not
literally ``cylinder area minus one pipe'' at all heights.

To make this precise, consider a horizontal slice at elevation $\eta$. Let
\begin{itemize}
  \item $n(\eta)$ be the number of pipe cross-sections intersected by that slice,
  \item $A$ be the cross-sectional area of a single pipe leg.
\end{itemize}
The effective water-bearing area in that slice is then
\[
  A_{\text{eff}}(\eta) = A_{\text{cyl}} - n(\eta)\,A,
\]
and the exact container water volume (excluding the pipe interior) is
\[
  V_{\text{cont}}(t)
  = \int_0^{h_c(t)} A_{\text{eff}}(\eta)\,d\eta
  = \int_0^{h_c(t)} \bigl[A_{\text{cyl}} - n(\eta)\,A\bigr]\,d\eta.
\]

During the filling phase the total water volume is conserved, so
\[
  V_{\text{cont}}(t) + V_{\text{pipe}}(t)
  = V_{\text{cont}}(0) + V_{\text{pipe}}(0).
\]
Introducing
\[
  F(h) = \int_0^{h} A_{\text{eff}}(\eta)\,d\eta,
\]
we can write
\[
  F(h_c(t)) + A L(t) = F(h_0) + A L_1.
\]
Equivalently, expressing the container level as a function of the wetted length,
\[
  F(h_c(L)) = F(h_0) + A (L_1 - L).
\]
In the fully general case, $h_c(L)$ is therefore obtained from an implicit relation involving
$F$ and its inverse:
\[
  h_c(L) = F^{-1}\Big(F(h_0) + A (L_1 - L)\Big).
\]

For a meandering pipe with straight legs, $A_{\text{eff}}(\eta)$ and hence $F(h)$ are piecewise
constant/linear functions of $h$, so $h_c(L)$ is piecewise linear with slope changes at
elevations where the number $n(\eta)$ of intersected legs changes.

In this paper, to keep the analytical model compact and highlight the main dynamical
mechanisms, we adopt the simpler relation
\[
  V_{\text{cont}}(t) = A_0 h_c(t),
\]
with $A_0$ treated as a constant. This corresponds to either:
\begin{enumerate}
  \item an idealised container geometry in which the pipe occupies a shaft of constant footprint
        across the range of $h_c$, or
  \item a slender-pipe approximation for the real geometry, in which the pipe diameter is small
        compared to the container radius and the variation of $n(\eta)$ with height only
        produces a small perturbation in the free-surface height.
\end{enumerate}
Under either interpretation, taking $A_0$ as constant leads to the linear volume--height
relation~\eqref{eq:global-volume} used below and greatly simplifies the coupling between the container level and the
advancing front, while retaining a faithful description of the dominant siphon dynamics.
For a given internal layout, the function $A_{\text{eff}}(\eta)$ and its integral
$F(h)=\int_0^{h}A_{\text{eff}}(\eta)\,d\eta$ provide an exact (though piecewise-defined)
map from container level to container volume; the constant-$A_0$ model is a compact
specialisation that captures the leading-order behaviour.

\section{Validity limits and missing physics}
\label{sec:limits}
The slug model developed below is intentionally minimal. It remains accurate only while
certain physical assumptions hold. Here we collect the main validity limits.
\begin{itemize}
  \item \textbf{Hydrostatic inlet assumption.} The inlet pressure is taken as
    $p(0,t)=p_{\mathrm{atm}}+\rho g\bigl(h_c(t)-h_0\bigr)$, i.e.\ the container is assumed to
    remain hydrostatic at each instant. This is valid when the free surface moves slowly,
    \[
      |\dot h_c|\ll \sqrt{gD},\qquad |\ddot h_c|\ll g,
    \]
    so that unsteady and dynamic effects at the inlet are small. If these conditions are
    violated, an unsteady Bernoulli term can be interpreted as an additional effective head
    of order
    \[
      \Delta h_{\text{unsteady}}\approx -\frac{A}{A_0}\,\frac{\dot h_c\,\dot L}{g},
    \]
    which lies beyond the present model.

  \item \textbf{Cavitation / tensile failure.} Along the centreline, and especially near the
    crest, the absolute pressure must remain above vapour pressure (or the tensile strength
    of the liquid) to avoid column separation. A conservative requirement is that the
    available head between container and crest, minus the loss head to that point, exceeds
    the vapour head,
    \begin{equation}
      h_c - h_{\text{crest}} - h_{\text{loss,crest}} \ge h_{\mathrm{vap}},\qquad
      h_{\mathrm{vap}} = \frac{p_{\mathrm{atm}}-p_{\mathrm{vap}}}{\rho g}.
      \label{eq:cavitation}
    \end{equation}
    If \eqref{eq:cavitation} is violated, the modelled negative pressures are not sustainable
    and the real flow is likely to cavitate or entrain air. Modern discussions of the siphon
    height limit and negative pressures include Boatwright et al.~\cite{boatwright} and
    Vera et al.~\cite{vera}; the role of water under tension in plant xylem is reviewed by
    Tyree~\cite{tyree}.

  \item \textbf{Capillarity / surface tension.} For very small diameters $D$ the meniscus
    at the inlet can sustain a finite pressure jump. A simple estimate is the Jurin-like
    height
    \[
      h_{\sigma}\approx \frac{4\sigma\cos\theta}{\rho g D},
    \]
    where $\sigma$ is the surface tension and $\theta$ the contact angle. A net head
    at least this large is required to break the meniscus and start the motion; once a finite
    velocity is established, the standard loss model applies.

  \item \textbf{Slug-model breakdown.} The present model treats the water as a single,
    one-dimensional slug with a sharp front. Strong three-dimensional curvature, vigorous
    turbulence, air entrainment, or repeated column separation can all invalidate this
    picture. In such regimes a two-phase or fully resolved CFD approach is more
    appropriate.
\end{itemize}

\section{Initial conditions at \texorpdfstring{$t=0$}{t=0}}

We choose $t = 0$ at the start of the filling phase, when the dynamical motion of the
water slug begins. In practice there are two simple wetting procedures that prepare this
state:
\begin{enumerate}
  \item \emph{Pre-fill to $B$.} The container is slowly filled until the free surface reaches
        point $B$ on the first leg, fully wetting segment $AB$. During this quasi-static
        process the wetted length simply tracks the container level; at the moment the
        surface passes $B$ one has $L=L_1$.
  \item \emph{Dip to $B$.} The inlet is dipped into a reservoir so that the water rises inside
        the leg until the front reaches $B$, after which the leg is placed in its operating
        position. Again, the filling up to $B$ is slow compared with the subsequent motion,
        so the front arrives at $B$ with negligible speed.
\end{enumerate}
In both cases the slow preparatory motion ends with the same configuration: the front is
at the end of segment $AB$, the container surface passes through $A$ at height $h_0$, and
the slug is momentarily at rest. It is therefore natural, and physically realistic, to adopt
the initial conditions
\begin{equation}
  t_0 = 0,\qquad h_c(0)=h_0,\qquad L(0) = L_1,\qquad \dot L(0) = 0.
  \label{eq:IC}
\end{equation}
The total water volume in the container + pipe system at $t = 0$ is
\[
  V_{\text{tot}}
  = V_{\text{cont}}(0) + V_{\text{pipe}}(0)
  = A_0 h_0 + A L_1.
\]
This quantity $V_{\text{tot}}$ will be crucial in the volume-based derivation of $h_c(L)$ during
the filling phase.

\section{Volume conservation in the filling phase}

\subsection{Global volume balance (container + pipe)}

During the filling phase, no water has yet left the outlet, so the total volume of water in the
combined container--pipe system remains constant:
\[
  V_{\text{tot}} = V_{\text{cont}}(t) + V_{\text{pipe}}(t) = \text{constant}.
\]
Using
\[
  V_{\text{cont}}(t) = A_0 h_c(t),\qquad
  V_{\text{pipe}}(t) = A L(t),
\]
we obtain
\[
  A_0 h_c(t) + A L(t) = V_{\text{tot}}.
\]
At $t = 0$ we already found $V_{\text{tot}} = A_0 h_0 + A L_1$. Therefore, for any later time
$t$ in the filling phase,
\begin{equation}
  A_0 h_c(t) + A L(t) = A_0 h_0 + A L_1.
  \label{eq:global-volume}
\end{equation}
Equation~\eqref{eq:global-volume} is the global volume conservation relation for the filling phase.

\subsection{Container level as a function of wetted length, $h_c(L)$}

Solving the conservation relation~\eqref{eq:global-volume} for $h_c$ in terms of $L$ gives
\[
  A_0 h_c(L) = A_0 h_0 + A L_1 - A L,
\]
so that
\[
  h_c(L) = h_0 - \frac{A}{A_0}\,(L - L_1).
\]
Thus during the filling phase,
\begin{equation}
  h_c(L) = h_0 - \frac{A}{A_0}(L - L_1).
  \label{eq:hc-of-L}
\end{equation}
This is the key relation: the container free-surface elevation expressed directly as a function
of the wetted length, obtained from plain volume conservation.

Important observations:
\begin{itemize}
  \item $h_c(L)$ decreases linearly as $L$ grows.
  \item The volume of water occupying any part of the pipe is already accounted for via the
        term $A L$. One does not need a separate ``pipe volume under the surface'' term;
        it is automatically included in $V_{\text{pipe}} = A L$.
  \item In the filling phase, $h_c(t)$ is not an independent state: it is fully determined by
        $L(t)$.
\end{itemize}

\section{Dynamics in the filling phase}

In the filling phase the pipe is only partially full: the front is at position $L(t)$ and the outlet
at $E$ is still dry. The water in the pipe is modelled as a one-dimensional slug of length $L(t)$,
cross-sectional area $A$, and uniform mean velocity
\[
  v(t) = \dot L(t).
\]
Its mass is
\[
  m(t) = \rho A L(t),
\]
where $\rho$ is the fluid density.

The driving force comes from the instantaneous head difference between the container
free surface at height $h_c(t)$ and the front elevation $h(L(t))$, while wall friction and local
disturbances in the pipe generate resistive forces. In this section we first derive the lossless
equation of motion from momentum conservation and then incorporate the loss models.

\subsection{Lossless momentum balance}

In the absence of friction and minor losses, the net driving force on the slug is
\[
  F_{\text{drive}}(t) = \rho g A \bigl[h_c(t) - h(L(t))\bigr],
\]
where $g$ is the gravitational acceleration. This result is derived below by integrating the
pressure and gravity contributions along the pipe.

Newton's second law in the direction of motion reads
\[
  m(t)\,\ddot L(t) = F_{\text{drive}}(t).
\]
Substituting $m(t) = \rho A L(t)$ and cancelling $\rho A$ yields
\[
  L(t)\,\ddot L(t) = g\bigl[h_c(t) - h(L(t))\bigr].
\]

During the filling phase the container level $h_c(t)$ is not an independent state variable:
it is determined by global volume conservation, and we may regard $h_c$ as a function of the wetted length via \eqref{eq:hc-of-L}.
Using this relation, the lossless filling-phase equation of motion can be written in closed form as
\begin{equation}
  \ddot L
  = \frac{g}{L}\Big[h_0 - \frac{A}{A_0}(L - L_1) - h(L)\Big].
  \label{eq:filling-lossless}
\end{equation}
Equation~\eqref{eq:filling-lossless} is integrated from the initial conditions
\[
  L(0) = L_1,\qquad \dot L(0) = v_0,
\]
up to the time $t^*$ at which the front first reaches the outlet, $L(t^*) = L_{\max}$.

\subsection{Detailed derivation of the net driving force}

We now show how the expression
\[
  F_{\text{drive}} = \rho g A \bigl[h_c(t) - h(L)\bigr]
\]
arises from pressure and gravity acting on the water in the pipe.

Introduce a local coordinate $\xi$ along the pipe centreline, measured from the inlet, so that
$\xi$ runs from $0$ at the entrance to $L(t)$ at the front. At position $\xi$ the pipe
centreline is at height $h(\xi)$, the pressure is $p(\xi,t)$, and the tangent to the pipe makes
some angle to the horizontal. Consider a thin fluid slice between $\xi$ and $\xi + d\xi$:
\begin{itemize}
  \item cross-sectional area: $A$,
  \item thickness: $d\xi$,
  \item volume: $A\,d\xi$,
  \item mass: $dm = \rho A\,d\xi$.
\end{itemize}
Forces tangential to the pipe on this slice are:
\begin{itemize}
  \item pressure force from the upstream side: $p(\xi,t)\,A$,
  \item pressure force from the downstream side: $-p(\xi + d\xi,t)\,A$,
  \item component of the weight along the pipe: $-\rho g A\,\frac{dh}{d\xi}\,d\xi$,
\end{itemize}
where the sign convention is chosen so that increasing $\xi$ follows the direction of motion
and $dh/d\xi$ accounts for the local inclination of the pipe. The total tangential force on the
slice is therefore
\[
  dF = -A\frac{\partial p}{\partial \xi}\,d\xi - \rho g A\,\frac{dh}{d\xi}\,d\xi.
\]
Integrating from $\xi = 0$ to $\xi = L(t)$ gives the net force on the whole slug:
\[
  F_{\text{net}}
  = \int_0^{L} dF
  = -A\int_0^{L}\frac{\partial p}{\partial \xi}\,d\xi
    - \rho g A \int_0^{L}\frac{dh}{d\xi}\,d\xi.
\]
The second integral is purely geometric:
\[
  \int_0^{L} \frac{dh}{d\xi}\,d\xi = h(L) - h(0),
\]
so
\[
  -\rho g A \int_0^{L} \frac{dh}{d\xi}\,d\xi
  = \rho g A\bigl[h(0) - h(L)\bigr].
\]
For the pressure term we obtain
\[
  -A\int_0^{L}\frac{\partial p}{\partial \xi}\,d\xi
  = -A\bigl[p(L,t) - p(0,t)\bigr]
  = A\bigl[p(0,t) - p(L,t)\bigr].
\]
Combining the two contributions yields
\[
  F_{\text{net}} = A\bigl[p(0,t) - p(L,t)\bigr]
                  + \rho g A\bigl[h(0) - h(L)\bigr].
\]

We now relate the end pressures $p(0,t)$ and $p(L,t)$ to the container free surface and the
outlet:
\begin{itemize}
  \item At the free surface in the container the pressure is $p_{\text{atm}}$. Hydrostatic
        equilibrium in the container gives the pressure at the pipe entrance,
        \[
          p(0,t) = p_{\text{atm}} + \rho g\bigl[h_c(t) - h(0)\bigr].
        \]
  \item At the outlet the jet discharges to the atmosphere at height $h(L)$, so
        \[
          p(L,t) = p_{\text{atm}}.
        \]
\end{itemize}
Substituting into the expression for $F_{\text{net}}$ and cancelling $p_{\text{atm}}$,
\[
  F_{\text{net}}
  = \rho g A\bigl[h_c(t) - h(0)\bigr]
    + \rho g A\bigl[h(0) - h(L)\bigr]
  = \rho g A\bigl[h_c(t) - h(L)\bigr].
\]
Thus the combined effect of the hydrostatic pressure in the container and the weight of
the water in the pipe is equivalent to a single driving head between the container free
surface and the front. In the lossless case $F_{\text{drive}} = F_{\text{net}}$, and inserting this
into Newton's law recovers the lossless equation of motion~\eqref{eq:filling-lossless}.

\subsection{Loss models: distributed friction and local (minor) losses}

In reality the pipe walls and geometric disturbances dissipate energy. We model these effects
using standard head-loss expressions.

\subsubsection*{Distributed wall friction (Darcy--Weisbach)}

The head loss due to viscous friction along a length $L$ of straight pipe with diameter $D$
and mean velocity $v$ is written in Darcy--Weisbach form as
\[
  h_f = f\,\frac{L}{D}\,\frac{v^2}{2g},
\]
where $f$ is the Darcy friction factor, which depends on the Reynolds number
\[
  \text{Re} = \frac{\rho v D}{\mu}
\]
and the relative roughness of the pipe. For laminar flow ($\text{Re} \lesssim 2000$) one has
\[
  f = \frac{64}{\text{Re}},
\]
while for turbulent and transitional flow an explicit all-regime correlation due to
Churchill~\cite{churchill} is particularly convenient. In the form given by
Churchill~\cite{churchill} it reads
\begin{align}
  \frac{f}{8} &= \left[\left(\frac{8}{\text{Re}}\right)^{12}
  + \frac{1}{\bigl(A_1 + A_2\bigr)^{3/2}}\right]^{1/12}, \label{eq:churchill-f} \\
  A_1 &= \left[-2.457 \ln\left(\left(\frac{7}{\text{Re}}\right)^{0.9}
  + 0.27\,\frac{\varepsilon}{D}\right)\right]^{16}, \label{eq:churchill-A1} \\
  A_2 &= \left(\frac{37530}{\text{Re}}\right)^{16}, \label{eq:churchill-A2}
\end{align}
where $\varepsilon/D$ is the relative roughness.
Equations~\eqref{eq:churchill-f}--\eqref{eq:churchill-A2} are used to evaluate the friction factor $f$ in the loss models below.
For comparison or validation one may still refer to the implicit
Colebrook--White equation or a Moody diagram~\cite{colebrook,moody}.
In many siphon applications a representative constant value of $f$ (based on the typical
operating $\text{Re}$) is adequate.
\subsubsection*{Local (minor) losses}

Bends, entrances, exits, contractions, valves, and other fittings are represented by local (or
``minor'') losses of the form
\[
  h_m = K\,\frac{v^2}{2g},
\]
where $K$ is a dimensionless loss coefficient determined by the fitting geometry. If several
elements are present along the wetted portion of the pipe, their effects are additive:
\[
  h_{\text{minor,tot}} = \Big(\sum_i K_i\Big)\,\frac{v^2}{2g}.
\]
Published tables give typical $K_i$ values for common components (sharp or rounded
entrances, elbows, valves, sudden expansions, and so on), and compilations such as
Crane~\cite{crane} are widely used in engineering practice.

\subsubsection*{Total head loss and resistive force}

Over a wetted length $L$ with mean velocity $v$ the total head loss is modelled as
\[
  h_{\text{loss}}(v,L)
  = f\,\frac{L}{D}\,\frac{v^2}{2g}
    + \Big(\sum_i K_i\Big)\,\frac{v^2}{2g}.
\]
In the slug model this head corresponds to a resistive force opposing the motion,
\[
  F_{\text{loss}} = \rho g A\,h_{\text{loss}}(v,L)\,\operatorname{sgn}(v).
\]
Using $v = \dot L$ in the filling phase and writing $|v|v = v\,\operatorname{sgn}(v)\,v$ gives
\[
  F_{\text{loss}} = \frac{\rho A}{2}\,\Big(f\,\frac{L}{D} + \sum_i K_i\Big)\,|\dot L|\,\dot L.
\]

\subsection{Extended equation of motion in the filling phase}

Including both the driving force and the resistive force, Newton's second law for the slug
becomes
\begin{equation}
  \rho A L\,\ddot L
  = \rho g A\bigl[h_c(t) - h(L)\bigr]
    - \frac{\rho A}{2}\,\Big(f\,\frac{L}{D} + \sum_i K_i\Big)\,|\dot L|\,\dot L.
  \label{eq:filling-extended-raw}
\end{equation}
Cancelling $\rho A$ and again using $h_c(L)$ from above, we obtain the extended filling-phase
ODE
\begin{equation}
  \ddot L
  = \frac{g}{L}\Big[h_0 - \frac{A}{A_0}(L - L_1) - h(L)\Big]
    - \frac{1}{2L}\,\Big(f\,\frac{L}{D} + \sum_i K_i\Big)\,|\dot L|\,\dot L.
  \label{eq:filling-extended}
\end{equation}
Equation~\eqref{eq:filling-extended} combines geometry, container drawdown, and both distributed and local
losses into a single nonlinear ODE for the front position $L(t)$ during the filling phase.
The initial conditions at $t = 0$ are
\[
  L(0) = L_1,\qquad \dot L(0) = v_0,
\]
with $h_c(L)$ given by the volume-conservation relation above. The filling phase ends when
$L(t)$ first reaches $L_{\max}$.

\subsection{Idealised lossless priming model and upper bound}

For design purposes it is useful to consider an idealised upper-bound model for the priming
process. In this model we make two additional simplifications:
\begin{itemize}
  \item the container level is held fixed at $h_c(t) \equiv h_0$ (no drawdown), and
  \item all losses are neglected, $f = 0$ and $\sum_i K_i = 0$.
\end{itemize}
We retain the same initial conditions as in the full model,
\[
  L(0) = L_1,\qquad \dot L(0) = 0,
\]
so that the front starts from rest at the end of the first segment.

Under these assumptions the filling-phase equation reduces to
\begin{equation}
  \ddot L = \frac{g}{L}\,\bigl[h_0 - h(L)\bigr].
  \label{eq:ideal-ddotL}
\end{equation}
Introducing $v(L) = \dot L$ and using $\ddot L = v\,dv/dL$ gives
\begin{equation}
  v\,\frac{dv}{dL}
  = \frac{g}{L}\,\bigl[h_0 - h(L)\bigr].
  \label{eq:ideal-vdv}
\end{equation}
Integrating from $L_1$ (where $v=0$) to a general $L$ yields
\begin{equation}
  \frac{1}{2}\bigl[v(L)^2 - 0\bigr]
  = g\int_{L_1}^{L} \frac{h_0 - h(\xi)}{\xi}\,d\xi,
  \label{eq:ideal-int}
\end{equation}
so that
\begin{equation}
  v(L)^2
  = 2g\int_{L_1}^{L} \frac{h_0 - h(\xi)}{\xi}\,d\xi.
  \label{eq:ideal-v2}
\end{equation}

\subsubsection*{Closed form for piecewise-linear $h(L)$}
On each geometric segment $i$ the elevation can be written in the affine form
\eqref{eq:linear-hi}, $h(L)=c_i+s_iL$ for $L\in [S_{i-1},S_i]$. The contribution of a
subinterval $[a,b]\subseteq [S_{i-1},S_i]$ to the integral in \eqref{eq:ideal-v2} is then
\begin{align}
  \int_{a}^{b} \frac{h_0 - h(\xi)}{\xi}\,d\xi
  &= \int_a^b \frac{h_0 - c_i - s_i\xi}{\xi}\,d\xi \notag\\
  &= (h_0-c_i)\int_a^b \frac{d\xi}{\xi} - s_i\int_a^b d\xi \notag\\
  &= (h_0-c_i)\ln\!\left(\frac{b}{a}\right) - s_i(b-a).
\end{align}
It is convenient to denote this segment contribution by
\begin{equation}
  \Phi_i(a,b) = (h_0-c_i)\ln\!\left(\frac{b}{a}\right) - s_i(b-a).
  \label{eq:Phi}
\end{equation}
For a given candidate $L$ lying on segment $k$, i.e.\ $L\in [S_{k-1},S_k]$, the integral in
\eqref{eq:ideal-v2} splits into full segments and a partial segment,
\begin{equation}
  \int_{L_1}^{L} \frac{h_0 - h(\xi)}{\xi}\,d\xi
  = \sum_{i=2}^{k-1} \Phi_i(S_{i-1},S_i) + \Phi_k(S_{k-1},L).
\end{equation}
The stall condition \eqref{eq:stall-condition} on that segment can thus be written as
\begin{equation}
  \sum_{i=2}^{k-1} \Phi_i(S_{i-1},S_i) + \Phi_k(S_{k-1},L) = 0,
  \label{eq:stall-piecewise}
\end{equation}
which is a monotone equation for $L$ on $[S_{k-1},S_k]$ and can be solved with a simple
root-finding method such as bisection. The resulting $L_{\text{stall}}$ gives the idealised
maximum climb $h_{\max}=h(L_{\text{stall}})$ for the chosen geometry.
In this ideal model the front comes to rest at the first arclength $L = L_{\text{stall}}$ such that
$v(L_{\text{stall}}) = 0$, i.e.
\begin{equation}
  \int_{L_1}^{L_{\text{stall}}} \frac{h_0 - h(\xi)}{\xi}\,d\xi = 0.
  \label{eq:stall-condition}
\end{equation}
The corresponding maximum centreline elevation achievable by the priming slug in the ideal
lossless model is
\begin{equation}
  h_{\max} = h\bigl(L_{\text{stall}}\bigr).
  \label{eq:hmax}
\end{equation}
Because this model ignores both drawdown and dissipation, $h_{\max}$ is an upper bound on
what can be attained in the full Churchill+loss model for the same geometry and initial
configuration.

\subsection{Geometric priming criterion}

Let $L_D$ denote the arclength coordinate of the crest point $D$; by construction its elevation is $h_3$.
The idealised lossless model above leads to a simple geometric necessary condition for
self-priming:
\begin{equation}
  h_{\max} > h_3.
\end{equation}
If $h_{\max} \le h_3$, then even the most optimistic model (no drawdown, no losses, front
starting from rest at $L_1$) predicts that the front stalls below the crest. In that case the
given geometry cannot self-prime, regardless of the detailed friction or wetting conditions.
If $h_{\max} > h_3$, the geometry passes this necessary condition; actual self-priming then
depends on the full dynamics with continuous Churchill friction and finite container
drawdown.


\section{Dynamics in the draining phase}

Once the front reaches the outlet, the pipe fills completely and the system transitions to a
draining phase. The wetted length is then fixed at
\[
  L(t) = L_{\max},\qquad t \ge t^*,
\]
where $t^*$ is the time of first filling, but the mean velocity $v(t)$ and the container level
$h_c(t)$ continue to evolve.

\subsection{Container mass balance}

In the draining phase the volumetric flow rate through the pipe is
\[
  Q(t) = A v(t),
\]
and this flow rate leaves the container. The container volume therefore satisfies
\[
  \frac{dV_{\text{cont}}}{dt} = -Q(t) = -A v(t).
\]
Using $V_{\text{cont}}(t) = A_0 h_c(t)$ and dividing by $A_0$ gives
\begin{equation}
  \dot h_c(t) = -\frac{A}{A_0}\,v(t),\qquad t \ge t^*.
  \label{eq:drain-hc}
\end{equation}
Unlike in the filling phase, the total water volume in container plus pipe is no longer
constant, so there is no algebraic relation of the form $h_c(L)$.

\subsection{Momentum balance with losses in the draining phase}

For a completely filled pipe the slug length is $L_{\max}$, the outlet elevation is the fixed height
\[
  h_{\text{out}} = h(L_{\max}),
\]
and the mean velocity $v(t)$ is the dynamical variable.

The momentum balance along the pipe now reads
\[
  \rho A L_{\max}\,\dot v
  = \rho g A\bigl[h_c(t) - h_{\text{out}}\bigr]
    - \rho g A\,h_{\text{loss}}(v,L_{\max})\,\operatorname{sgn}(v),
\]
with
\[
  h_{\text{loss}}(v,L_{\max})
  = \Big(f\,\frac{L_{\max}}{D} + \sum_i K_i\Big)\,\frac{v^2}{2g}.
\]
Cancelling $\rho A$ yields the draining-phase velocity equation
\begin{equation}
  L_{\max}\,\dot v
  = g\bigl[h_c(t) - h_{\text{out}}\bigr]
    - \frac{1}{2}\,\Big(f\,\frac{L_{\max}}{D} + \sum_i K_i\Big)\,|v|\,v.
  \label{eq:drain-v-L}
\end{equation}
Equivalently, we may write
\begin{equation}
  \dot v
  = \frac{g}{L_{\max}}\bigl[h_c - h_{\text{out}}\bigr]
    - \frac{1}{2L_{\max}}\,\Big(f\,\frac{L_{\max}}{D} + \sum_i K_i\Big)\,|v|\,v.
  \label{eq:drain-v}
\end{equation}
Together, the equations~\eqref{eq:drain-hc} and~\eqref{eq:drain-v}
form a coupled system for $h_c(t)$ and $v(t)$ in the draining phase.

The initial conditions at the start of draining are taken from the end of the filling phase:
\[
  v(t^*) = \dot L(t^*),\qquad
  h_c(t^*) = h_0 - \frac{A}{A_0}(L_{\max} - L_1).
\]

\subsection{Quasi-steady draining approximation}

In many practical situations the acceleration term $L_{\max}\,\dot v$ is small compared with the
friction and minor losses, especially once the flow has developed. Neglecting inertia by
setting $\dot v \approx 0$ in the velocity equation~\eqref{eq:drain-v} gives a quasi-steady relation between $v$
and $h_c$:
\[
  0 \approx \frac{g}{L_{\max}}\bigl[h_c(t) - h_{\text{out}}\bigr]
            - \frac{1}{2L_{\max}}\,\Big(f\,\frac{L_{\max}}{D} + \sum_i K_i\Big)\,v\,|v|.
\]
Restricting attention to the usual draining regime with $v(t) \ge 0$ and $h_c(t) \ge h_{\text{out}}$,
we obtain
\begin{equation}
  v(t) \approx
  \sqrt{\frac{2g\bigl[h_c(t) - h_{\text{out}}\bigr]}
             {f\,L_{\max}/D + \sum_i K_i}}.
  \label{eq:quasi-steady-v}
\end{equation}
Substituting~\eqref{eq:quasi-steady-v} into the mass-balance equation~\eqref{eq:drain-hc} leads to a single first-order ODE for
$h_c(t)$ that reproduces the familiar textbook formula for gravity-driven drainage with
friction and minor losses. The full unsteady system, however, is required to describe the
transient immediately after priming.

\section{Phase transition and overall summary}

The complete model for a self-starting siphon based on the present geometry and volume
description can be summarised as follows; it links the centreline geometry $h(L)$ introduced in \cref{sec:geometry} to the dynamical variables $L(t)$, $v(t)$, and $h_c(t)$.
\begin{enumerate}
  \item \textbf{Filling phase ($0 \le t < t^*$).}
    \begin{itemize}
      \item State variables: wetted length $L(t)$ and its time derivative $\dot L(t)$.
      \item Container level is constrained by volume conservation:
            \[
              h_c(L) = h_0 - \frac{A}{A_0}(L - L_1).
            \]
      \item Dynamics: the front position $L(t)$ evolves according to the nonlinear ODE
            given in \eqref{eq:filling-extended},
            with initial data $L(0) = L_1$ and $\dot L(0) = 0$.
      \item The filling phase ends at the first time $t^*$ such that $L(t^*) = L_{\max}$.
    \end{itemize}
  \item \textbf{Draining phase ($t \ge t^*$).}
    \begin{itemize}
      \item State variables: mean pipe velocity $v(t)$ and container level $h_c(t)$.
      \item Kinematics: the pipe remains full ($L = L_{\max}$) and the container height obeys
            the mass-balance equation \eqref{eq:drain-hc}.
      \item Dynamics: the velocity satisfies the momentum equation \eqref{eq:drain-v},
            with initial conditions
            $v(t^*) = \dot L(t^*)$ and
            $h_c(t^*) = h_0 - \frac{A}{A_0}(L_{\max} - L_1)$.
    \end{itemize}
  \item \textbf{Loss models.}
    \begin{itemize}
      \item Distributed wall friction is modelled via the Darcy--Weisbach relation with
            friction factor $f$ determined from Reynolds number and pipe roughness.
      \item Local geometric disturbances are represented by minor-loss coefficients $K_i$.
      \item In both phases, distributed and local losses enter the dynamics through the
            combination
            \[
              f\,\frac{L}{D} + \sum_i K_i,
            \]
            evaluated with $L = L(t)$ in the filling phase and $L = L_{\max}$ in the
            draining phase.
    \end{itemize}
\end{enumerate}
These equations provide a unified, self-consistent description of the priming and draining
dynamics for a self-starting siphon, explicitly connecting the geometry $h(L)$, container
drawdown $h_c(L)$, and frictional and minor losses to the evolution of $L(t)$, $v(t)$, and
$h_c(t)$ in time.

\subsection*{Variables and parameters}
For quick reference, the main symbols appearing in the model are:
\begin{description}
  \item[$L(t)$] wetted length of the slug along the pipe centreline, measured from $A$ to the front.
  \item[$h(L)$] centreline elevation at arclength $L$; piecewise linear according to the geometry.
  \item[$h_c(t)$] container free-surface elevation measured from the container base.
  \item[$h_0$] initial container height, $h_c(0)$, and vertical coordinate of point $A$.
  \item[$v(t)$] mean slug velocity, $v=\dot L$, in the filling phase; mean pipe velocity in the draining phase.
  \item[$A$] pipe cross-sectional area (circular diameter $D$), $A=\pi D^2/4$.
  \item[$A_0$] effective container cross-sectional area in the constant-area model.
  \item[$L_i$] length of segment $i$ of the pipe; $L_{\max}=L_1+L_2+L_3+L_4$.
  \item[$\rho,\mu$] fluid density and dynamic viscosity.
  \item[$f$] Darcy friction factor, evaluated from the Churchill correlation \eqref{eq:churchill-f}--\eqref{eq:churchill-A2}.
  \item[$\sum_i K_i$] total minor-loss coefficient over the wetted portion of the pipe.
\end{description}

\section{Computational approach}
\label{sec:comp}
The coupled ODEs derived above are well suited to numerical integration. A convenient
strategy is to treat the filling and draining phases as two separate initial-value problems
linked by an event.
\begin{itemize}
  \item \textbf{Filling phase.} Integrate the second-order ODE \eqref{eq:filling-extended}
        for $L(t)$ (or, equivalently, the first-order system for $(L,\dot L)$) starting from
        the initial conditions \eqref{eq:IC}. At each step, evaluate the container height from
        $h_c(L)=h_0-(A/A_0)(L-L_1)$ and update the friction factor $f$ using the instantaneous
        Reynolds number. Terminate the integration at the event $L(t^*)=L_{\max}$, recording
        $t^*$ and the front speed $v(t^*)=\dot L(t^*)$.
  \item \textbf{Draining phase.} Use $h_c(t^*) = h_0 - \tfrac{A}{A_0}(L_{\max}-L_1)$ and
        $v(t^*)$ as initial conditions for the draining system \eqref{eq:drain-hc}--\eqref{eq:drain-v}.
        Integrate forward until either $h_c(t)$ approaches $h_{\text{out}}$ or the velocity
        drops to zero.
  \item \textbf{Loss evaluation.} In both phases, compute the combination
        $fL/D+\sum_i K_i$ using the current wetted length $L$ (or $L_{\max}$ in draining)
        and update $f$ from the Churchill formula as the Reynolds number evolves.
\end{itemize}
Standard ODE solvers such as Runge--Kutta methods with adaptive step size (e.g.\ those
described in Hairer et al.~\cite{hairer}) are adequate. In pseudocode, a simple implementation
in a high-level language might look like:
\begin{verbatim}
integrate filling_rhs from t=0 with state (L=L1, v=0),
  stopping when L = Lmax;
record t_star, v_star, and h_c_star;
integrate drain_rhs from t=t_star with state (h_c=h_c_star, v=v_star),
  stopping when h_c = h_out or v = 0.
\end{verbatim}
This approach mirrors the physical phase change from a partially filled, advancing slug
to a fully filled draining siphon and allows direct comparison between simulation and
experiment for a given geometry.

\begin{thebibliography}{99}

\bibitem{white}
F.~M.~White,
\emph{Fluid Mechanics}.
McGraw--Hill.


\bibitem{churchill}
S.~W.~Churchill,
``Friction-factor equation spans all fluid-flow regimes,''
\emph{Chemical Engineering}, Vol.~84, No.~24, pp.~91--92 (1977).


\bibitem{colebrook}
C.~F.~Colebrook,
``Turbulent flow in pipes, with particular reference to the transition region
between smooth and rough pipe laws,''
Journal of the Institution of Civil Engineers (1939).

\bibitem{moody}
L.~F.~Moody,
``Friction factors for pipe flow,''
Transactions of the ASME (1944).

\bibitem{crane}
Crane Co.,
\emph{Flow of Fluids Through Valves, Fittings, and Pipe (Technical Paper 410)}.
Crane Co.

\bibitem{potter_barnes}
A.~Potter and F.~H.~Barnes,
``The siphon,''
Physics Education \textbf{6}, 362--366 (1971).

\bibitem{jumper_stanchev}
W.~D.~Jumper and B.~Stanchev,
``Towards explaining the water siphon,''
The Physics Teacher \textbf{52}, 474--478 (2014).

\bibitem{zheng_siphon}
J.~Zheng, J.~Guo, J.~Wang, Y.~Zhang, Q.~Lu, and H.~Sun,
``Calculation of the flow velocity of a siphon,''
Physics of Fluids \textbf{33}, 017105 (2021).

\bibitem{ortloff_kassinos}
C.~R.~Ortloff and A.~Kassinos,
``Computational fluid dynamics investigation of the hydraulic behaviour of the Roman inverted siphon system at Aspendos, Turkey,''
Journal of Archaeological Science \textbf{30}(4), 417--428 (2003).

\bibitem{sedlar_siphon}
M.~Sedlar, P.~Prochazka, M.~Komarek, V.~Uruba, and V.~Skala,
``Experimental research and numerical analysis of flow phenomena in discharge object with siphon,''
Water \textbf{12}, 3330 (2020).

\bibitem{hairer}
E.~Hairer, S.~P.~Norsett, and G.~Wanner,
\emph{Solving Ordinary Differential Equations I: Nonstiff Problems}, 2nd ed.
Springer, Berlin, 1993.

\bibitem{boatwright}
A.~Boatwright, S.~Hughes, and J.~Barry,
``The height limit of a siphon,''
Scientific Reports \textbf{5}, 16790 (2015).

\bibitem{vera}
F.~Vera, R.~Rivera, D.~Romero-Maltrana, and J.~Villanueva,
``Negative pressures and the first water siphon taller than 10.33 meters,''
PLOS ONE \textbf{11}(4), e0153055 (2016).

\bibitem{tyree}
M.~T.~Tyree,
``Plant hydraulics: The ascent of water,''
Nature \textbf{423}, 923 (2003).

\end{thebibliography}

\end{document}
